\chapter{État de l'Art et Étude Théorique}
\label{chap:etat_art}

\section{Introduction}

L'application développée pour CIVCO BTP s'appuie sur un ensemble de technologies web
modernes, réparties entre un \textit{backend} API, une base de données relationnelle et
une interface utilisateur riche. Ce chapitre constitue le socle théorique du rapport :
il présente les concepts et les technologies étudiés, puis justifie les choix retenus
pour la réalisation du projet.

\medskip

Nous abordons successivement les \textbf{architectures des applications web}, le
\textbf{framework Laravel} côté serveur, les \textbf{interfaces React en SPA}, les
\textbf{bases de données relationnelles et l'ORM}, les \textbf{API REST et
l'authentification}, ainsi que les \textbf{solutions ERP et outils de gestion de projets
dans le secteur du BTP}. Le chapitre se conclut par une référence comparative
(\textit{benchmark}) justifiant chaque choix technologique retenu.

% ─────────────────────────────────────────────────────────────────────
\section{Les applications web et l'architecture full-stack}
\subsection{Évolution des applications web}

Les applications de gestion d'entreprise ont connu une évolution majeure au cours des
deux dernières décennies. Les premières solutions reposaient sur des architectures
\textbf{monolithiques} : logique métier, interface et base de données regroupées dans
un même programme, souvent accessible via un navigateur avec rechargement complet de
page à chaque action.

\medskip

L'émergence des \textbf{API} (\textit{Application Programming Interface}) a séparé la
couche de présentation de la couche métier. Le navigateur ou une application mobile
consomme des services distants via \textbf{HTTP}, échangeant des données structurées en
\textbf{JSON}. Cette découpe a ouvert la voie aux architectures \textit{full-stack}
modernes, dans lesquelles un \textit{frontend} JavaScript riche dialogue avec un
\textit{backend} API indépendant~\cite{fielding2000rest}.

\subsection{Architecture client-serveur et tiers}

L'architecture \textbf{three-tier} (\textit{trois tiers}) structure l'application en
trois couches distinctes :

\begin{enumerate}
  \item \textbf{Couche présentation} : interface utilisateur (React), exécutée dans le
  navigateur.
  \item \textbf{Couche métier} : logique applicative, règles de gestion, authentification
  (Laravel).
  \item \textbf{Couche données} : persistance relationnelle (MySQL / PostgreSQL).
\end{enumerate}

Cette séparation améliore la maintenabilité : chaque couche peut évoluer indépendamment
tant que les contrats d'interface (API REST) sont respectés. Fowler~\cite{fowler2019patterns}
décrit ce découpage comme un patron fondamental des applications d'entreprise.

\begin{figure}[H]
  \centering
  \includegraphics[width=0.55\textwidth]{fig_2_1_three_tier_preview.png}
  \caption{Architecture three-tier de l'application CIVCO BTP}
  \label{fig:three_tier}
\end{figure}

\subsection{Single Page Application (SPA)}

Une \textbf{SPA} (\textit{Single Page Application}) charge une seule page HTML, puis
met à jour dynamiquement le contenu via JavaScript sans rechargement complet. Les
avantages pour une application métier sont significatifs :

\begin{itemize}
  \item \textbf{Fluidité} : navigation instantanée entre les modules (projets, clients,
  factures).
  \item \textbf{Expérience utilisateur} : interfaces réactives, proches des applications
  de bureau.
  \item \textbf{Réutilisation de l'état} : conservation du contexte utilisateur
  (filtres, session) entre les écrans.
\end{itemize}

Le couplage d'une SPA React avec une API Laravel constitue aujourd'hui un standard
industriel pour les tableaux de bord et outils de gestion internes.

\subsection{Panorama des stacks full-stack}

\begin{table}[H]
\caption{Panorama des stacks full-stack pour applications métier}
\label{tab:stack_panorama}
\centering
\small
\begin{tabular}{|l|l|p{8cm}|}
\hline
\textbf{Stack} & \textbf{Langages} & \textbf{Caractéristiques} \\
\hline
LAMP / LEMP & PHP + MySQL & Mature, hébergement abordable, Laravel/Symfony \\
\hline
MERN / MEAN & JavaScript & Un seul langage, Node.js + MongoDB \\
\hline
Django + React & Python + JS & Rapidité backend, ORM Django \\
\hline
\textbf{Laravel + React} & PHP + JS & API REST + SPA, écosystème riche \\
\hline
.NET + Angular & C\# + TS & Écosystème Microsoft, entreprise \\
\hline
\end{tabular}
\end{table}

Pour CIVCO BTP, la stack \textbf{Laravel + React} (Tableau~\ref{tab:stack_panorama})
offre le meilleur compromis entre productivité, courbe d'apprentissage et adéquation
avec les compétences visées en filière IIR : développement web, bases de données et
architecture logicielle.

\subsection{Échange de données et format JSON}

Les API modernes échangent des données structurées en \textbf{JSON} (\textit{JavaScript
Object Notation}), format léger et lisible, natif en JavaScript et bien supporté par
PHP/Laravel. Chaque ressource métier est sérialisée en objet JSON : un projet expose
son identifiant, son client associé, ses phases, son statut et ses dates ; un devis
expose ses lignes, montants et statut de validation.

\medskip

Ce format facilite le découplage frontend/backend : React consomme directement les
réponses JSON, les mappe vers ses composants et met à jour l'interface sans
transformation complexe. Les codes de statut HTTP (\texttt{200}, \texttt{201},
\texttt{422}, \texttt{403}) signalent le résultat de chaque opération de façon
standardisée.

\subsection{Patterns CRUD et ressources métier}

La majorité des modules de l'application suivent le pattern \textbf{CRUD}
(\textit{Create, Read, Update, Delete}) : créer un client, lister les projets,
modifier une facture, archiver un document. Laravel formalise ce pattern via les
contrôleurs de ressources et les routes RESTful, ce qui garantit une cohérence
d'API sur l'ensemble des modules (clients, projets, devis, factures, tâches,
contrats).


% ─────────────────────────────────────────────────────────────────────
\section{Le framework Laravel et l'écosystème PHP}
\subsection{Présentation de PHP et Laravel}

\textbf{PHP} est un langage de script côté serveur largement déployé sur le web.
Historiquement associé au développement de sites dynamiques, il a évolué vers des
versions modernes (8.x) intégrant un typage strict, des performances accrues (JIT) et
un écosystème mature de frameworks.

\medskip

\textbf{Laravel}~\cite{laravel} est le framework PHP le plus adopté pour construire des
applications web et des API. Créé par Taylor Otwell, il impose une structure claire
autour du patron \textbf{MVC} (\textit{Model-View-Controller}) et propose des
composants intégrés pour les tâches courantes : routage, validation, authentification,
files d'attente, envoi de courriels et gestion des fichiers.

\subsection{Architecture MVC dans Laravel}

Le patron \textbf{MVC} répartit les responsabilités :

\begin{itemize}
  \item \textbf{Modèle} : représentation des entités métier et accès aux données
  (Eloquent).
  \item \textbf{Vue} : dans une API, les \textit{Resources} et \textit{Collections}
  formatent les réponses JSON pour le client.
  \item \textbf{Contrôleur} : reçoit la requête HTTP, applique la validation, invoque
  les services métier et retourne la réponse.
\end{itemize}

Cette organisation facilite la lecture du code et l'ajout de nouveaux modules. Pour
l'application CIVCO BTP, chaque domaine métier (projets, clients, devis, factures)
dispose de son contrôleur, de ses modèles et de ses règles de validation.

\subsection{Composants clés de l'écosystème Laravel}

\subsubsection{Routage et validation}

Le fichier de routes (\texttt{routes/api.php}) déclare les points d'entrée de l'API.
Chaque route est associée à un contrôleur et protégée par des \textit{middleware}
(vérification d'authentification, contrôle de permissions, contexte entreprise).

Les \textbf{Form Requests} encapsulent les règles de validation des données entrantes
(ex.~: champs obligatoires d'un devis, formats de dates, montants positifs), garantissant
l'intégrité des données avant persistance.

\subsubsection{Services et logique métier}

Au-delà des contrôleurs, la logique complexe est extraite dans des classes
\textbf{Service} (\texttt{QuoteService}, \texttt{PrintTrackingService}, etc.). Ce
découpage respecte le principe de responsabilité unique : le contrôleur orchestre,
le service implémente la règle métier.

\subsubsection{Migrations et seeders}

Les \textbf{migrations} versionnent le schéma de la base de données sous forme de fichiers
PHP exécutables de façon reproductible. Les \textbf{seeders} initialisent des données
de référence (rôles, permissions, données de démonstration), indispensables pour
tester l'application en conditions proches du réel.

\subsection{Comparaison avec les alternatives PHP}

\begin{table}[H]
\caption{Comparaison des frameworks PHP pour une API REST}
\label{tab:php_frameworks}
\centering
\small
\begin{tabular}{|l|p{4cm}|p{4cm}|p{4cm}|}
\hline
\textbf{Framework} & \textbf{Philosophie} & \textbf{Forces} & \textbf{Contexte idéal} \\
\hline
\textbf{Laravel} & Convention over configuration & Écosystème complet, Sanctum, Eloquent &
API + SPA, PME/ETI \\
\hline
Symfony & Composants réutilisables & Modularité, robustesse & Grandes applications \\
\hline
CodeIgniter & Minimaliste & Légèreté & Petits projets \\
\hline
\end{tabular}
\end{table}

Laravel ressort du Tableau~\ref{tab:php_frameworks} pour un projet de gestion métier
nécessitant authentification, RBAC, documents et évolutivité --- le profil exact du
projet CIVCO BTP.

\subsection{Laravel dans le projet CIVCO BTP}

L'API Laravel expose plus de quarante contrôleurs organisés par domaine, des
\textit{middleware} de contrôle d'accès et un ensemble de services métier. Cette
structure permet d'ajouter progressivement de nouveaux modules (avenants contractuels,
modèles de documents, suivi d'impression) sans remettre en cause l'architecture
existante.

\subsection{Middleware, Resources et couche de présentation API}

\subsubsection{Middleware}

Les \textbf{middleware} Laravel interceptent chaque requête HTTP avant qu'elle
n'atteigne le contrôleur. Dans le projet, plusieurs middlewares assurent la
sécurité et le contexte métier :

\begin{itemize}
  \item \textbf{Authentification Sanctum} : vérifie que l'utilisateur est connecté.
  \item \textbf{Contrôle de permissions} : valide que l'utilisateur possède la
  permission requise (\texttt{project.view}, \texttt{quote.manage}, etc.).
  \item \textbf{Contexte entreprise} : identifie la société courante via l'en-tête
  \texttt{X-Company-Id}, garantissant l'isolation des données métier.
\end{itemize}

\subsubsection{API Resources}

Les classes \textbf{Resource} (\texttt{ProjectResource}, \texttt{InvoiceResource},
etc.) transforment les modèles Eloquent en tableaux JSON normalisés. Elles permettent
de contrôler précisément les champs exposés à l'interface, d'inclure conditionnellement
des relations (\texttt{client}, \texttt{lines}, \texttt{teamMembers}) et d'uniformiser
le format de réponse sur toute l'API.

\subsubsection{Validation et gestion des erreurs}

Lorsqu'une requête échoue la validation, Laravel retourne une réponse \texttt{422}
avec le détail des champs en erreur. Le frontend React intercepte ces réponses pour
afficher des messages à l'utilisateur, améliorant l'expérience de saisie sur les
formulaires de devis, factures et fiches projets.


% ─────────────────────────────────────────────────────────────────────
\section{React et les Single Page Applications}
\subsection{Définition et principes des SPA}

Une \textbf{SPA} (\textit{Single Page Application}) est une application web dont
l'interface est construite dynamiquement en JavaScript. Contrairement aux sites
classiques où chaque clic déclenche un rechargement HTML complet, la SPA charge le
shell de l'application une fois, puis met à jour le DOM (\textit{Document Object Model})
au fil des interactions utilisateur.

\medskip

Ce modèle repose sur le \textbf{routing côté client} : la bibliothèque React Router
associe chaque URL à un composant React (\texttt{/projects}, \texttt{/clients},
\texttt{/invoices}), sans solliciter le serveur pour le HTML. Seules les données
nécessaires sont récupérées via l'API REST.

\subsection{React : bibliothèque declarative}

\textbf{React}~\cite{react}, développé par Meta, est une bibliothèque JavaScript pour
construire des interfaces utilisateur par \textbf{composants} réutilisables. Chaque
composant encapsule structure (JSX), état et comportement.

\medskip

Les principes fondamentaux de React sont :

\begin{itemize}
  \item \textbf{Composants fonctionnels} : fonctions JavaScript retournant du JSX,
  composées hiérarchiquement (Sidebar, Dashboard, ProjectDetailPage, etc.).
  \item \textbf{État local} (\texttt{useState}) : données modifiables au sein d'un
  composant (formulaires, filtres, modales).
  \item \textbf{Effets de bord} (\texttt{useEffect}) : chargement de données API au
  montage du composant ou lors d'un changement de paramètre.
  \item \textbf{Context API} : partage d'état global (authentification, thème, langue)
  sans propager les props à travers toute l'arborescence.
\end{itemize}

Cette approche declarative simplifie la construction d'interfaces complexes : l'état
de l'interface est une fonction de l'état applicatif.

\subsection{Écosystème React du projet}

L'application CIVCO BTP s'appuie sur plusieurs bibliothèques complémentaires :

\begin{table}[H]
\caption{Composants de l'écosystème frontend}
\label{tab:react_ecosystem}
\centering
\small
\begin{tabular}{|l|l|p{7.5cm}|}
\hline
\textbf{Outil} & \textbf{Rôle} & \textbf{Usage dans le projet} \\
\hline
Vite~\cite{vite} & Build tool & Compilation rapide, HMR en développement \\
\hline
React Router & Routing & Navigation entre modules ERP \\
\hline
Axios & Client HTTP & Appels API REST avec intercepteurs \\
\hline
Tailwind CSS~\cite{tailwindcss} & Styles utilitaires & Interface sombre cohérente \\
\hline
Recharts & Graphiques & Tableau de bord, KPIs \\
\hline
Leaflet & Cartographie & Carte des projets / chantiers \\
\hline
\end{tabular}
\end{table}

\subsection{Communication avec le backend}

Le \textit{frontend} React consomme l'API Laravel via des modules dédiés
(\texttt{api/projects.js}, \texttt{api/clients.js}, etc.). Un client Axios centralisé
configure :

\begin{itemize}
  \item l'envoi automatique des cookies de session (Sanctum) ;
  \item le jeton CSRF pour les requêtes \texttt{POST}, \texttt{PUT} et \texttt{DELETE} ;
  \item la gestion des erreurs \texttt{401} (redirection vers la page de connexion).
\end{itemize}

Ce patron \textit{API-first} garantit que toute action utilisateur (créer un devis,
archiver un client, assigner une tâche) transite par la couche métier Laravel, où les
règles de validation et de permissions sont appliquées.

\begin{figure}[H]
\centering
\resizebox{0.92\textwidth}{!}{%
\begin{tikzpicture}[
  box/.style={rectangle, rounded corners=4pt, fill=ensamBlue!12, draw=ensamBlue,
              minimum width=3.2cm, minimum height=1.0cm, align=center,
              font=\small\bfseries},
  arr/.style={-{Stealth[length=6pt]}, thick, draw=gray!70},
  lbl/.style={font=\footnotesize, text=gray!70}
]
  \node[box] (browser) at (0, 0)   {Navigateur\\React SPA};
  \node[box] (api)     at (5.5, 0) {API Laravel\\REST + Sanctum};
  \node[box] (db)      at (11, 0)  {Base de\\données};

  \draw[arr] (browser) -- (api) node[midway, above, lbl] {JSON / HTTPS};
  \draw[arr] (api)     -- (db)  node[midway, above, lbl] {Eloquent / SQL};
\end{tikzpicture}
}%
\caption{Flux de communication React $\rightarrow$ Laravel $\rightarrow$ SGBD}
\label{fig:react_flow}
\end{figure}

\subsection{Comparaison React / Vue / Angular}

\begin{table}[H]
\caption{Comparaison des principaux frameworks frontend (2025-2026)}
\label{tab:spa_compare}
\centering
\small
\begin{tabular}{|l|p{3.5cm}|p{3.5cm}|p{3.5cm}|}
\hline
\textbf{Critère} & \textbf{React} & \textbf{Vue.js} & \textbf{Angular} \\
\hline
Flexibilité & Élevée (bibliothèque) & Élevée (progressif) & Opinionated (framework complet) \\
\hline
Courbe d'apprentissage & Modérée & Faible & Élevée \\
\hline
Marché de l'emploi & Très large & Large & Large (entreprise) \\
\hline
Écosystème SPA & React Router, Vite & Vue Router, Vite & Router intégré \\
\hline
\end{tabular}
\end{table}

React est retenu pour sa flexibilité, son adoption industrielle et la disponibilité de
bibliothèques adaptées aux tableaux de bord métier (graphiques, cartes, composants
modulaires), comme le montre le Tableau~\ref{tab:spa_compare}.

\subsection{Organisation du code frontend}

Le code React est structuré par \textbf{features} (\texttt{features/projects/},
\texttt{features/clients/}, \texttt{features/invoices/}), chacune regroupant pages,
composants et utilitaires propres au module. Les composants transverses (Sidebar,
PermissionGate, modales) vivent dans \texttt{components/}. Cette organisation par
domaine facilite la navigation dans le code et l'ajout de nouvelles fonctionnalités
sans impacter les modules existants.

\subsection{Interface utilisateur et Tailwind CSS}

L'interface adopte un design sombre professionnel, cohérent sur tous les modules.
\textbf{Tailwind CSS}~\cite{tailwindcss} est un framework CSS utilitaire : les classes
(\texttt{flex}, \texttt{rounded-lg}, \texttt{bg-gray-800}) s'appliquent directement
dans le JSX. Cette approche accélère le prototypage et garantit une cohérence visuelle
sur les écrans de l'application (listes, fiches détail, modales, tableaux de bord).

\medskip

Le \textbf{tableau de bord} illustre l'intérêt de React pour la visualisation métier :
graphiques (Recharts), widgets configurables, indicateurs \textbf{KPI} et accès rapide
aux projets en cours. La carte des chantiers (Leaflet) offre une vue géographique des
projets, pertinente pour une entreprise disposant de plusieurs sites simultanés.

\subsection{Internationalisation}

L'application supporte le \textbf{français} (langue par défaut) et l'\textbf{anglais}
via un contexte de langue et des fichiers de traduction (\texttt{locales/fr.js},
\texttt{locales/en.js}). Cette préparation facilite l'usage de l'outil par des
collaborateurs anglophones ou une future extension à d'autres contextes linguistiques.


% ─────────────────────────────────────────────────────────────────────
\section{Les bases de données relationnelles et l'ORM}

\subsection{Modèle relationnel et SGBD}

Une application de gestion de projets manipule des entités fortement structurées et
liées entre elles : clients, projets, phases, tâches, devis, factures, contrats,
utilisateurs et rôles. Le \textbf{modèle relationnel}, formalisé par Codd, organise ces
données en tables reliées par des clés étrangères, garantissant l'intégrité référentielle
et la cohérence des enregistrements~\cite{fowler2019patterns}.

\medskip

Les \textbf{SGBD} (\textit{Systèmes de Gestion de Bases de Données}) relationnels les
plus répandus en environnement web sont \textbf{MySQL}~\cite{mysql} et
\textbf{PostgreSQL}~\cite{postgresql}. MySQL se distingue par sa simplicité de déploiement
et sa large adoption ; PostgreSQL offre des fonctionnalités avancées (contraintes
complexes, types JSON, performances sur de gros volumes). Pour une application métier
comme celle de CIVCO BTP, les deux systèmes conviennent ; le choix final dépend de
l'environnement d'hébergement et des besoins de montée en charge.

\subsection{L'ORM et Eloquent}

Manipuler directement le langage \textbf{SQL} dans chaque composant applicatif alourdit
le code et complique la maintenance. L'\textbf{ORM} (\textit{Object-Relational Mapping})
fait correspondre les tables de la base à des classes du langage de programmation.
Dans Laravel, \textbf{Eloquent} fournit une couche d'abstraction élégante :

\begin{itemize}
  \item chaque table est représentée par un \textbf{modèle} (\texttt{Project},
  \texttt{Client}, \texttt{Invoice}, etc.) ;
  \item les relations (\texttt{hasMany}, \texttt{belongsTo}, \texttt{belongsToMany})
  expriment les liens métier de façon déclarative ;
  \item les migrations versionnent le schéma de la base et facilitent le déploiement
  entre environnements ;
  \item les \textit{seeders} permettent de peupler la base avec des données de démonstration.
\end{itemize}

Cette approche accélère le développement, réduit les erreurs de requêtes et aligne le
code objet sur le modèle métier de l'entreprise.

\subsection{Normalisation et intégrité des données}

La conception d'une base pour un ERP léger repose sur plusieurs principes :

\begin{itemize}
  \item \textbf{Normalisation} : éviter la duplication d'informations (un client est
  enregistré une seule fois, référencé par ses projets et documents).
  \item \textbf{Contraintes d'intégrité} : clés étrangères, unicité des identifiants
  métier (numéros de devis, de facture).
  \item \textbf{Traçabilité} : horodatage des enregistrements (\texttt{created\_at},
  \texttt{updated\_at}) et journalisation des actions sensibles.
\end{itemize}

Ces principes sont essentiels dans un contexte BTP où les documents commerciaux et
contractuels doivent rester cohérents et auditables.

\subsection{Modèle de données métier}

Le schéma relationnel de l'application modélise les entités principales du cycle de
vie d'un projet CIVCO BTP :

\begin{table}[H]
\centering
\small
\caption{Entités principales du modèle de données}
\label{tab:entities}
\begin{tabular}{|l|p{9cm}|}
\hline
\textbf{Entité} & \textbf{Rôle} \\
\hline
Client & Maître d'ouvrage ou client commercial, contacts associés \\
\hline
Project & Chantier / projet, lié à un client, phases et équipe \\
\hline
ProjectPhase & Découpage temporel du projet (études, exécution, réception) \\
\hline
Task & Tâche opérationnelle assignée à un collaborateur \\
\hline
Quote / Invoice & Documents commerciaux avec lignes de détail \\
\hline
Contract & Contrat et avenants liés au projet \\
\hline
User / Role & Utilisateurs internes et permissions RBAC \\
\hline
\end{tabular}
\end{table}

Les relations entre ces entités (un client possède plusieurs projets, un projet
contient plusieurs phases et tâches, un devis peut être converti en facture) sont
implémentées via Eloquent et contraintes de clés étrangères en base.

% ─────────────────────────────────────────────────────────────────────
\section{Les API REST et l'authentification}

\subsection{Principe des API REST}

Une \textbf{API REST} (\textit{Representational State Transfer}) expose des ressources
(identifiées par des URL) manipulables via les verbes \textbf{HTTP} :
\texttt{GET} (lecture), \texttt{POST} (création), \texttt{PUT/PATCH} (mise à jour),
\texttt{DELETE} (suppression)~\cite{fielding2000rest}. Ce style architectural sépare
clairement le \textit{frontend} (interface utilisateur) du \textit{backend} (logique
métier et persistance), ce qui permet :

\begin{itemize}
  \item de faire évoluer l'interface sans réécrire le serveur ;
  \item de réutiliser la même API pour d'éventuels clients mobiles futurs ;
  \item de tester indépendamment chaque couche.
\end{itemize}

Dans le projet CIVCO BTP, l'API est organisée sous le préfixe \texttt{/api/v1/} et
expose des ressources métier : projets, clients, devis, factures, tâches, etc.

\subsection{Authentification et contrôle d'accès}

La sécurisation d'une application métier repose sur deux piliers :

\medskip

\textbf{Authentification.} Laravel \textbf{Sanctum}~\cite{sanctum} gère l'authentification
par session pour les applications \textbf{SPA} : le navigateur conserve un cookie de
session sécurisé, complété par une protection \textbf{CSRF} (\textit{Cross-Site Request
Forgery}) sur les requêtes mutantes. Cette approche est adaptée à une application hébergée
sur le même domaine que l'API, sans la complexité d'un flux JWT (\textit{JSON Web Token})
côté client.

\medskip

\textbf{Autorisation (RBAC).} Le contrôle d'accès basé sur les rôles (\textbf{RBAC})
associe à chaque utilisateur un ensemble de permissions (\texttt{project.view},
\texttt{quote.manage}, \texttt{invoice.view}, etc.). Les routes API vérifient ces
permissions avant d'exécuter l'action demandée. Ainsi, un commercial, un chef de
projet et un administrateur accèdent chacun aux modules correspondant à leur profil
métier.

% ─────────────────────────────────────────────────────────────────────
\section{ERP et digitalisation du secteur BTP}

\subsection{Enjeux de la gestion intégrée en BTP}

Le secteur du \textbf{BTP} partage avec l'industrie des besoins de traçabilité, de
planification et de coordination documentaire. Un \textbf{ERP} (\textit{Enterprise
Resource Planning}) vise à unifier ces processus : relation client, suivi de chantier,
achats, facturation, ressources humaines. Dans le BTP, les spécificités suivantes
renforcent l'intérêt d'un outil intégré :

\begin{itemize}
  \item multiplicité des intervenants (MOA, BET, entreprises, sous-traitants) ;
  \item cycle documentaire riche (devis, bons de commande, BL, factures, avenants) ;
  \item suivi temporel des phases et des tâches de chantier ;
  \item besoin de visibilité consolidée pour la direction et les conducteurs de travaux.
\end{itemize}

\subsection{Panorama des solutions existantes}

Le marché propose des solutions generalistes et des outils spécialisés construction :

\begin{table}[H]
\centering
\small
\caption{Panorama de solutions de gestion applicables au BTP}
\label{tab:erp_panorama}
\begin{tabular}{|l|l|p{7.5cm}|}
\hline
\textbf{Solution} & \textbf{Type} & \textbf{Apport principal} \\
\hline
SAP S/4HANA & ERP generaliste & Suite intégrée grande entreprise, modules finance et projets \\
\hline
Sage 100 / X3 & ERP PME & Comptabilité, commercial, gestion de production \\
\hline
Procore~\cite{procore} & SaaS BTP & Gestion de projets construction, documents, coordination \\
\hline
Buildertrend & SaaS BTP & Suivi chantier, relation client, devis \\
\hline
Solution sur mesure & Développement interne & Adaptation exacte aux processus de l'entreprise \\
\hline
\end{tabular}
\end{table}

Les ERP generalistes offrent une couverture fonctionnelle large mais imposent des coûts
de licence, de paramétrage et de formation élevés. Les solutions SaaS BTP sont plus
ciblées mais restent standardisées. Pour CIVCO BTP, le développement d'une
\textbf{application sur mesure} permet d'aligner précisément l'outil sur les processus
internes, tout en maîtrisant l'évolution fonctionnelle.

\subsection{Digitalisation des BET au Maroc}

Les bureaux d'études techniques marocains, comme CIVCO BTP, combinent expertise
ingénierie et activités de suivi de chantier. La digitalisation de leurs processus
internes --- gestion documentaire, suivi des échéances, coordination inter-services ---
constitue un levier de compétitivité. Une application web interne permet de :

\begin{itemize}
  \item réduire le temps de recherche d'information sur un projet ;
  \item standardiser la production des devis et factures ;
  \item améliorer la visibilité de l'avancement pour la direction et les conducteurs ;
  \item préparer une traçabilité des actions (journal d'activité, historique).
\end{itemize}

Ce contexte sectoriel justifie l'investissement dans une solution adaptée plutôt que
dans un ERP generaliste surdimensionné pour la taille et les processus de l'entreprise.

% ─────────────────────────────────────────────────────────────────────
\section{Références comparatives (Benchmarks) et justification des choix}

\subsection{Comparaison des frameworks backend}

\begin{table}[H]
\centering
\small
\caption{Comparaison des frameworks backend pour une API métier}
\label{tab:backend_benchmark}
\begin{tabular}{|l|p{3.2cm}|p{3.2cm}|p{3cm}|}
\hline
\textbf{Framework} & \textbf{Points forts} & \textbf{Limites} & \textbf{Écosystème} \\
\hline
\textbf{Laravel} & ORM Eloquent, migrations, Sanctum, rapidité de dev & Performance brute inférieure à Go/Rust & Très riche (PHP) \\
\hline
Symfony & Modularité, robustesse entreprise & Courbe d'apprentissage plus longue & Large (PHP) \\
\hline
Django & Admin intégré, ORM puissant & Moins adapté aux SPA modernes & Python \\
\hline
Express.js & Légèreté, JavaScript full-stack & Structure laissée au développeur & Node.js \\
\hline
\end{tabular}
\end{table}

\textbf{Laravel}~\cite{laravel} ressort du Tableau~\ref{tab:backend_benchmark} pour sa
productivité, son ORM mature, son système d'authentification Sanctum adapté aux SPA et
sa documentation exhaustive --- critères décisifs pour un projet de stage avec délai
contraint.

\subsection{Comparaison des frameworks frontend}

\begin{table}[H]
\centering
\small
\caption{Comparaison des frameworks frontend}
\label{tab:frontend_benchmark}
\begin{tabular}{|l|c|c|c|c|}
\hline
\textbf{Framework} & \textbf{Popularité} & \textbf{Écosystème} & \textbf{Courbe} &
\textbf{SPA} \\
\hline
\textbf{React} & Très élevée & Très riche & Modérée & Native \\
\hline
Vue.js & Élevée & Riche & Faible & Native \\
\hline
Angular & Élevée & Complet (Google) & Élevée & Native \\
\hline
Svelte & Croissante & Moyen & Faible & Native \\
\hline
\end{tabular}
\end{table}

\textbf{React}~\cite{react} est retenu (Tableau~\ref{tab:frontend_benchmark}) pour son
écosystème (React Router, bibliothèques UI), sa adoption industrielle massive et sa
compatibilité avec \textbf{Vite}~\cite{vite} pour un cycle de développement rapide.

\subsection{Comparaison des SGBD}

\begin{table}[H]
\centering
\small
\caption{Comparaison des systèmes de gestion de bases de données}
\label{tab:db_benchmark}
\begin{tabular}{|l|c|c|c|c|}
\hline
\textbf{SGBD} & \textbf{Open source} & \textbf{JSON} & \textbf{Performance} &
\textbf{Déploiement} \\
\hline
\textbf{MySQL} & Oui & Partiel & Élevée & Simple \\
\hline
\textbf{PostgreSQL} & Oui & Natif & Très élevée & Standard \\
\hline
SQLite & Oui & Oui & Faible (fichier) & Très simple \\
\hline
SQL Server & Non & Oui & Élevée & Entreprise \\
\hline
\end{tabular}
\end{table}

MySQL et PostgreSQL conviennent toutes deux au projet ; PostgreSQL est privilégié en
production pour sa robustesse, MySQL ou SQLite pouvant servir en développement local.

\subsection{Synthèse des choix}

\begin{table}[H]
\caption{Synthèse des choix technologiques et critères décisifs}
\label{tab:tech_choices}
\centering
\small
\begin{tabular}{|l|l|p{7cm}|}
\hline
\textbf{Composant} & \textbf{Choix retenu} & \textbf{Critère décisif} \\
\hline
Backend & Laravel 13 / PHP 8.3 & Productivité, ORM, écosystème mature \\
\hline
Frontend & React 19 + Vite 8 & SPA réactive, écosystème riche \\
\hline
Style UI & Tailwind CSS 4~\cite{tailwindcss} & Développement rapide, cohérence visuelle \\
\hline
API & REST / JSON & Simplicité, séparation frontend/backend \\
\hline
Authentification & Laravel Sanctum & Sessions SPA sécurisées, CSRF \\
\hline
Base de données & MySQL / PostgreSQL & Modèle relationnel, intégrité \\
\hline
ORM & Eloquent & Mapping objet-relationnel, migrations \\
\hline
Contrôle d'accès & RBAC (rôles/permissions) & Profils métier CIVCO BTP \\
\hline
\end{tabular}
\end{table}

Le Tableau~\ref{tab:tech_choices} récapitule l'ensemble des choix retenus pour
l'application de gestion de projets de CIVCO BTP.

% ─────────────────────────────────────────────────────────────────────
\section{Méthodologies de développement logiciel}

\subsection{Approches traditionnelles et agiles}

Le développement d'une application métier peut suivre différentes méthodologies.
L'approche \textbf{cascade} (\textit{Waterfall}) enchaîne les phases de manière
séquentielle : analyse, conception, développement, tests, déploiement. Elle convient
aux projets dont les besoins sont stables et parfaitement définis en amont, mais
peut s'avérer rigide lorsque les exigences évoluent en cours de réalisation.

\medskip

Les méthodes \textbf{agiles}~\cite{agilemanifesto}, à l'inverse, privilégient des
itérations courtes (sprints de deux à quatre semaines), une collaboration continue
avec le client et l'adaptation au changement. Le \textbf{Manifeste Agile} (2001)
formalise quatre valeurs : individus et interactions plutôt que processus et outils,
logiciel fonctionnel plutôt que documentation exhaustive, collaboration client plutôt
que négociation contractuelle, adaptation au changement plutôt que suivi d'un plan
rigide.

\subsection{Application au projet CIVCO BTP}

Dans le contexte de ce stage, une approche \textbf{itérative et incrémentale} a été
retenue, s'inspirant des principes agiles. Le suivi des tâches a été matérialisé via
\textbf{Trello}, organisé selon un tableau \textbf{Kanban} (colonnes \textit{À faire},
\textit{En cours}, \textit{Terminé}), sans appliquer un cadre Scrum formel (sprints,
daily meetings). Les raisons de ce choix sont les suivantes :

\begin{itemize}
  \item \textbf{Périmètre évolutif} : les besoins métier ont été affinés au fil des
  démonstrations, notamment pour les modules d'impression et de portail client.
  \item \textbf{Durée limitée} : quatre mois de stage imposent des livrables réguliers
  plutôt qu'une phase d'analyse prolongée.
  \item \textbf{Proximité métier} : l'encadrant professionnel était disponible pour
  valider chaque incrément et orienter les priorités.
  \item \textbf{Risque maîtrisé} : chaque module livré est testable indépendamment,
  réduisant le risque d'une intégration finale décevante.
\end{itemize}

\subsection{Qualité logicielle et bonnes pratiques}

Indépendamment de la méthodologie, plusieurs principes de qualité logicielle ont
guidé le développement :

\begin{itemize}
  \item \textbf{Principe DRY} (\textit{Don't Repeat Yourself}) : factorisation de
  la logique commune dans des services et des composants réutilisables.
  \item \textbf{Principe SOLID} : responsabilité unique des classes (un service par
  domaine métier), injection de dépendances via le conteneur Laravel.
  \item \textbf{Conventions de code} : Laravel Pint pour le PHP, ESLint pour le
  JavaScript, respect des standards PSR-12 et des conventions React.
  \item \textbf{Tests automatisés} : couverture des flux critiques (authentification,
  CRUD, conversion devis) par PHPUnit pour détecter les régressions.
  \item \textbf{Documentation} : commentaires sur les règles métier complexes,
  présent rapport et annexes techniques.
\end{itemize}

Ces pratiques, combinées à l'architecture modulaire décrite aux chapitres suivants,
visent à produire un code maintenable et évolutif, au-delà de la seule durée du stage.

\newpage
\section*{Conclusion}

Ce chapitre a posé les fondements théoriques et technologiques sur lesquels repose
l'application développée pour CIVCO BTP. Nous avons d'abord rappelé l'évolution des
applications web vers des architectures \textit{full-stack} en trois tiers, dans
lesquelles la couche présentation (React), la couche métier (Laravel) et la couche
données (SGBD relationnel) sont clairement séparées. Ce découplage facilite la
maintenance, les tests et l'évolution indépendante de chaque composant.

\medskip

L'étude de \textbf{Laravel} a montré comment un framework PHP mature structure le
\textit{backend} autour du patron \textbf{MVC}, d'Eloquent, de migrations et de
Sanctum. Les middleware, les Form Requests et les API Resources garantissent à la
fois la sécurité, la validation des données et un format de réponse JSON cohérent
pour l'ensemble des modules métier.

\medskip

Côté interface, \textbf{React} permet de construire une \textbf{SPA} modulaire et
réactive, organisée par fonctionnalités et appuyée sur un écosystème riche (Vite,
Axios, Tailwind CSS, Recharts, Leaflet). L'échange permanent entre le navigateur et
l'API REST via JSON constitue le lien central de l'architecture retenue.

\medskip

Le modèle relationnel et l'\textbf{ORM} assurent l'intégrité et la cohérence des
données --- clients, projets, phases, tâches, devis, factures et contrats --- tandis
que le \textbf{RBAC} adapte l'accès aux différents profils métier de l'entreprise.
Enfin, le panorama des solutions ERP et BTP a situé le projet dans son contexte
sectoriel et confirmé la pertinence d'une application sur mesure pour CIVCO BTP.

\medskip

Les benchmarks comparatifs ont permis de justifier chaque choix : Laravel pour la
productivité backend, React pour l'expérience utilisateur, PostgreSQL/MySQL pour la
persistance, Sanctum pour l'authentification SPA. L'ensemble de ces décisions est
guidé par trois impératifs : la \textbf{productivité} de développement dans le cadre
du stage, la \textbf{maintenabilité} du code sur le long terme, et l'\textbf{adéquation
métier} aux processus internes de l'entreprise.

\medskip

Le chapitre suivant traduit ces choix en une architecture système concrète : analyse
des besoins, identification des acteurs, cas d'utilisation et modélisation
\textbf{UML} de l'application de gestion et de planification des projets de
construction.

\newpage
